% ============ 第2章 §2.15 端点与暴露点(原书 pp.161-165) ============
\section{端点与暴露点}\label{sec:2.15}
\index{端点}
\index{暴露点}


在 $H^{p}$($1 \leq p \leq \infty$)的单位球 $S^{p}$ 中,作为 $\phi \in (H^{p})^{*}$ 的极值问题除常数倍外唯一解的 $f \in S^{p}$ 称为露点(exposed point).$S^{p}$ 的露点全部属于 $S^{p}$ 的边界.由引理~\ref{lem:2.14.3} 与定理~\ref{thm:2.14.4},当 $1 < p < \infty$ 时,$S^{p}$ 的边界上的函数全部是露点.p = 1,$\infty$ 时,露点是 $S^{p}$ 中非常特殊的函数.在 $S^{p}$ 中,不是 $S^{p}$ 的两个不同函数的真正凸组合的 $f \in S^{p}$ 称为端点(extreme point).$S^{p}$ 的端点全部属于 $S^{p}$ 的边界.当 $1 < p < \infty$ 时逆成立是 Minkowski 不等式的结果.本节考察 p = 1,$\infty$ 情形下的端点与露点.容易证明露点总是端点.在 $1 < p < \infty$ 时 $S^{p}$ 的露点与端点同 $S^{p}$ 的边界一致,但我们将证明在 p = 1,$\infty$ 时存在是端点但不是露点的边界点.

\begin{lemma}\label{lem:2.15.1}
设 p = 1 或 $\infty$.f 不是 $S^{p}$ 的端点的充要条件是,存在非零的 $g \in S^{p}$,使得 $\|f \pm g\|_{p} \leq 1$.
\end{lemma}

\begin{proof}
若 f 不是端点,则存在 0 < a < 1 与 $h, k \in S^{p}$($h \not\equiv k$),使得 $f = ah + (1-a)k$.取 $\varepsilon$ 使 $0 < a \pm \varepsilon < 1$.令 $\ell = f + \varepsilon(h - k)$ 且 $m = f - \varepsilon(h - k)$,则 $\ell$ 与 $m$ 属于 $S^{p}$.令 $g = (\ell - m)/2$,则 $g \not\equiv 0$,$f + g = \ell$ 且 $f - g = m$.逆命题由 $f = (f + g)/2 + (f - g)/2$ 显然.这个证明对 $1 < p < \infty$ 也成立.
\end{proof}

\begin{lemma}\label{lem:2.15.2}
设 $g, f \in H^{1}$ 且 a.e. t 有 $|g(e^{it})| \leq |f(e^{it})|$.若 f 是外函数,则 $|g(z)| \leq |f(z)|$($z \in D$).
\end{lemma}

\begin{proof}
由定理~\ref{thm:2.8.1},
\begin{equation*}
\log|g(re^{i\theta})| \leq \frac{1}{2\pi}\int_{-\pi}^{\pi} P(r, \theta - t)\log|g(e^{it})|\dd t,
\end{equation*}
又因 f 是外函数,
\begin{equation*}
\log|f(re^{i\theta})| = \frac{1}{2\pi}\int_{-\pi}^{\pi} P(r, \theta - t)\log|f(e^{it})|\dd t.
\end{equation*}
这证明了 $|g(re^{i\theta})| \leq |f(re^{i\theta})|$($0 \leq r < 1, -\pi \leq \theta < \pi$).
\end{proof}

\begin{theorem}\label{thm:2.15.3}
设 $f \in H^{1}$.f 是 $S^{1}$ 的端点的充要条件是 f 满足 $\|f\|_{1} = 1$ 且为外函数.
\end{theorem}

\begin{proof}
设 $f \in S^{1}$ 为外函数.若 f 不是端点,由引理~\ref{lem:2.15.1},存在非零 $g \in S^{1}$,使得 $\|f \pm g\|_{1} \leq 1$.由 $\|f\|_{1} = 1$ 得 $\|f + g\|_{1} = \|f - g\|_{1} = 1$.于是令 k = g/f,
\begin{equation*}
\int_{-\pi}^{\pi} \left\{|1 + k(e^{i\theta})| + |1 - k(e^{i\theta})| - 2\right\} |f(e^{i\theta})|\dd \theta \leq 0.
\end{equation*}
因 a.e. $\theta$ 有 $|1 + k(e^{i\theta})| + |1 - k(e^{i\theta})| \geq 2$,故 a.e. $\theta$ 有 $|1 + k(e^{i\theta})| + |1 - k(e^{i\theta})| = 2$,从而 a.e. $\theta$ 有 $-1 \leq k(e^{i\theta}) \leq 1$.于是 a.e. $\theta$ 有 $|g(e^{i\theta})| \leq |f(e^{i\theta})|$,由引理~\ref{lem:2.15.2},k 属于 $H^{\infty}$.由定理~\ref{thm:2.10.5},$\hat{k}(n) = 0$($n = \pm 1, \pm 2, \ldots$),由 Fourier 系数的唯一性,k 为常数,但由 $|1 + k(e^{i\theta})| = |1 - k(e^{i\theta})|$ 得 $k \equiv 0$.于是 $g \equiv 0$.这与假设矛盾,故 f 是端点.

设 f 为 $S^{1}$ 的端点.此时 $\|f\|_{1} = 1$.若 f 不是外函数,由定理~\ref{thm:2.8.3},f = qh 可作因式分解,其中 q 是非常值的内函数,h 是外函数.对实常数 $\alpha$,
\begin{equation*}
g = \frac{1}{2}(e^{i\alpha}q + e^{-i\alpha}\bar{q})f,
\end{equation*}
则 $g \in H^{1}$.因 q 不是常数,$g \not\equiv 0$.选取 $\alpha$ 使得
\begin{equation*}
\int_{-\pi}^{\pi} \mathrm{Re}\left\{e^{i\alpha}q(e^{i\theta})\right\} |f(e^{i\theta})|\dd \theta = 0,
\end{equation*}
则
\begin{equation*}
\|f \pm g\|_{1} = \int_{-\pi}^{\pi} |1 \pm \operatorname{Re}\left\{e^{i\alpha}q(e^{i\theta})\right\}| |f(e^{i\theta})|\dd \theta = 1.
\end{equation*}
这与引理~\ref{lem:2.15.1} 矛盾,故 f 是外函数.
\end{proof}

\begin{theorem}\label{thm:2.15.4}
设 $f \in H^{\infty}$.f 是 $S^{\infty}$ 的端点的充要条件是 $\|f\|_{\infty} = 1$ 且
\begin{equation*}
\int_{-\pi}^{\pi} \log(1 - |f(e^{i\theta})|)\dd \theta = -\infty.
\end{equation*}
\end{theorem}

\begin{proof}
设 $f \in S^{\infty}$ 满足
\begin{equation*}
\int_{-\pi}^{\pi} \log(1 - |f(e^{i\theta})|)\dd \theta = -\infty.
\end{equation*}
设 $g \in S^{\infty}$ 满足 $\|f \pm g\|_{\infty} \leq 1$,则
\begin{equation*}
|g(e^{i\theta})|^{2} \leq 1 - |f(e^{i\theta})|^{2} \leq 2(1 - |f(e^{i\theta})|),
\end{equation*}
故 $\int_{-\pi}^{\pi} \log|g(e^{i\theta})|\dd \theta = -\infty$.由定理~\ref{thm:2.5.3},$g \equiv 0$.于是由引理~\ref{lem:2.15.1},f 是 $S^{\infty}$ 的端点.

反之,若 $\|f\|_{\infty} = 1$ 且
\begin{equation*}
\int_{-\pi}^{\pi} \log(1 - |f(e^{i\theta})|)\dd \theta > -\infty,
\end{equation*}
令
\begin{equation*}
g(z) = \exp\left\{\frac{1}{2\pi}\int_{-\pi}^{\pi} \frac{e^{it} + z}{e^{it} - z}\log(1 - |f(e^{it})|)\dd t\right\},
\end{equation*}
则因 $\log(1 - |f(e^{it})|) \leq 0$,$|g(z)| \leq 1$,且由定理~\ref{thm:2.3.2},a.e. t 有 $|g(e^{it})| = 1 - |f(e^{it})|$.于是 $\|f \pm g\|_{\infty} \leq 1$ 且 $g \in S^{\infty}$,由引理~\ref{lem:2.15.1},f 不是 $S^{\infty}$ 的端点.
\end{proof}

\begin{proposition}\label{prop:2.15.5}
设 $f \in S^{1}$ 且 $\|f\|_{1} = 1$,则 f 是 $S^{1}$ 的露点的充要条件是:满足 $g \in H^{1}$ 且 a.e. $\theta$ 有 $g(e^{i\theta})/f(e^{i\theta}) \geq 0$ 的 g 仅为 f 的标量倍.
\end{proposition}

\begin{proof}
设 f 是 $S^{1}$ 的露点.可设存在 $\phi \in (H^{1})^{*}$ 且 $\phi(f) > 0$.若非零 $g \in H^{1}$ 满足 a.e. $\theta$ 有 $g(e^{i\theta})/f(e^{i\theta}) \geq 0$,则由引理~\ref{lem:2.14.3},$g/\|g\|_{1}$ 也是 $\phi$ 的极值函数.因 f 是 $S^{1}$ 的露点,g 是 f 的标量倍.设 f 不是 $S^{1}$ 的露点.令 K = $|f|/zf$,设 $\phi$ 为由 K 导出的 $(H^{1})^{*}$ 的元素,则 f 是 $\phi$ 的极值函数.因 f 不是露点,存在 $g \in S^{1}$,不是 f 的标量倍,且 $\phi(g) = \phi(f)$.由引理~\ref{lem:2.14.3},a.e. $\theta$ 有 $g(e^{i\theta})/f(e^{i\theta}) \geq 0$.
\end{proof}

\begin{proposition}\label{prop:2.15.6}
设 $f \in S^{\infty}$,若 f 在 $\partial D$ 上正测度的集合上绝对值等于 1,则 f 是 $S^{\infty}$ 的露点.
\end{proposition}

\begin{proof}
令 $E = \{\theta \in [-\pi, \pi]; |f(e^{i\theta})| = 1\}$,则 $\dd \theta(E) > 0$.
\begin{equation*}
K(e^{i\theta}) = \left\{ \begin{array}{c l}
e^{-i\theta}|f(e^{i\theta})|/f(e^{i\theta}) & (\theta \in E), \\
0 & (\theta \notin E)
\end{array} \right.
\end{equation*}
则 $K \in L^{1}$,且 a.e. $\theta$ 有 $e^{i\theta}f(e^{i\theta})K(e^{i\theta}) \geq 0$.由引理~\ref{lem:2.14.3},f 是由极值核 K 定义的 $\phi \in (H^{\infty})^{*}$ 的极值函数,且 $\phi(f) > 0$.设 g 为 $\phi$ 的另一极值函数,则存在实常数 $\alpha$,a.e. $\theta$ 有 $e^{i\alpha}e^{i\theta}g(e^{i\theta})K(e^{i\theta}) \geq 0$,故 $g = e^{-i\alpha}f$.
\end{proof}

\begin{example}\label{exa:2.15.7}
设 $f \in S^{1}$ 且 $\|f\|_{1} = 1$.若存在 $k \in H^{\infty}$($k^{-1} \in H^{\infty}$)使得 a.e. $\theta$ 有 $\operatorname{Re} f(e^{i\theta})k(e^{i\theta}) \geq 0$,则 f 是 $S^{1}$ 的端点.
\end{example}

\begin{proof}
首先证明当 $g = fk$ 时 g 是外函数.由定理~\ref{thm:2.8.3},$g = B_{1}S_{1}h_{1}$ 可作因式分解.因 a.e. $\theta$ 有 $\operatorname{Re} g(e^{i\theta}) \geq 0$,由定理~\ref{thm:2.9.1},$\operatorname{Re} g(z) > 0$($z \in D$),故 $g = S_{1}h_{1}$.因 $\operatorname{Re} g^{-1}(z) > 0$($z \in D$),由定理~\ref{thm:2.11.2},$g^{-1} \in \bigcap_{p<1} H^{p}$.由定理~\ref{thm:2.8.3},$g^{-1} = S_{2}h_{2}$ 可作因式分解.于是 $1 = S_{1}S_{2}h_{1}h_{2}$.由定义,对 j = 1,2,
\begin{equation*}
S_{j}(z) = e^{i\gamma_{j}}\exp\left\{-\int_{-\pi}^{\pi} \frac{e^{it} + z}{e^{it} - z}\dd \mu_{j}(t)\right\},
\end{equation*}
其中 $\mu_{j} \in M[-\pi, \pi]$ 是奇异正测度.又由定义,对 j = 1,2,
\begin{equation*}
h_{j}(z) = \exp\left\{\int_{-\pi}^{\pi} \frac{e^{it} + z}{e^{it} - z}\log|h_{j}(e^{it})|\dd t\right\}.
\end{equation*}
于是由 $1 = S_{1}S_{2}h_{1}h_{2}$,a.e. $\theta$ 有 $|(h_{1}h_{2})(e^{i\theta})| = 1$,故 $(h_{1}h_{2})(z) = 1$($z \in D$).于是 $(S_{1}S_{2})(z) = 1$($z \in D$),故 $\mu_{1} + \mu_{2} = 0$,从而 $\mu_{1} = \mu_{2} = 0$.于是 $S_{1} = e^{i\gamma_{1}}$,$g = e^{i\gamma_{1}}h_{1}$ 是外函数.k 是外函数可同样证明,故 $f = k^{-1}g$ 也是外函数.
\end{proof}

\begin{example}\label{exa:2.15.8}
设 $f \in S^{1}$ 且 $\|f\|_{1} = 1$.若存在 $k \in H^{\infty}$($k^{-1} \in H^{\infty}$)使得 a.e. $\theta$ 有 $\operatorname{Re} f(e^{i\theta})k(e^{i\theta}) \geq 0$,则 f 是 $S^{1}$ 的露点.
\end{example}

\begin{proof}
令 $g = fk/\|fk\|_{1}$.因 $k, k^{-1} \in H^{\infty}$,由命题~\ref{prop:2.15.5},只需证明:若 $h \in S^{1}$ 满足 a.e. $\theta$ 有 $h(e^{i\theta})/g(e^{i\theta}) \geq 0$,则 h = g.由 a.e. $\theta$ 有 $\mathrm{Re}(h(e^{i\theta}) + g(e^{i\theta}))/2 \geq 0$,由例~\ref{exa:2.15.7},$(h + g)/2$ 是外函数.另一方面,若 h 不是 g 的标量倍,则与定理~\ref{thm:2.15.3} 矛盾.
\end{proof}

定理~\ref{thm:2.15.3} 是 de Leeuw--Rudin 的结果 \cite{ch2:5}.已知命题~\ref{prop:2.15.6} 的逆也成立 \cite{ch2:6}.例~\ref{exa:2.15.8} 出自 \cite{ch2:7},其证明依据 \cite{ch2:8}.$S^{1}$ 的露点与端点相比研究得还不充分.关于露点的研究见 \cite{ch2:9} 与 \cite{ch2:10}.
